\section{Safety Compliance and International Standards}

Safety is a mandatory, non-negotiable requirement for industrial deployment; adherence to international safety standards acts as a facility's license to operate. The proposed dual-arm collaborative cell must strictly align with the frameworks established by the International Organization for Standardization (ISO).

\subsection{ISO 10218 and ISO/TS 15066 Integration}
Traditional industrial robots (such as the ABB IRB120) are governed by ISO 10218-1 and ISO 10218-2, which dictate safety requirements for industrial environments \cite{iso15066}. However, the inclusion of a collaborative workspace requires adherence to ISO/TS 15066. This technical specification provides comprehensive risk assessment guidance, including pain onset level studies used to establish maximum force and pressure thresholds for human-robot contact across 29 specific body areas \cite{robotiq2016iso}.

\subsection{Speed and Separation Monitoring (SSM)}
The primary safety protocol for collaborative overlap in this cell is Speed and Separation Monitoring (SSM) \cite{nist2016ssm}. ISO/TS 15066 mandates that a minimum protective distance be continuously maintained between the operator and the robotic system. The SSM equation calculates this protective distance based on variables including the human's travel speed, the robot's current speed, the system's response time, and the distance the robot travels during its braking sequence \cite{robotiq2016iso}. If the separation distance falls below the calculated threshold, the system must trigger a protective stop or seamlessly transition to a safe alternative path.

The integration of the \texttt{ZoneDetectionManager} in this project serves as the software-level precursor to this requirement, continuously evaluating spatial proximities to prevent both robot-robot and robot-human interference.